\textbf{Wärmetransport im Körper.} Die Wärme entsteht vor allem in den Muskeln und Organen (Stoffwechsel). Das Blut nimmt sie auf und verteilt sie im ganzen Körper (\emph{Konvektion}). Durch Gewebe und Haut gelangt sie durch \emph{Wärmeleitung} an die Oberfläche. Der Körper regelt den Transport: Bei Wärme werden die Blutgefässe der Haut weit gestellt, sodass mehr warmes Blut an die Oberfläche fliesst; bei Kälte werden sie verengt, und das Fettgewebe isoliert.

\textbf{Wärmeabgabe an die Umgebung.}
\begin{itemize}
\item \emph{Wärmestrahlung} (Infrarot) der Haut; in Ruhe bei Zimmertemperatur der grösste Anteil.
\item \emph{Konvektion}: Die Luft an der Haut wird erwärmt und steigt auf oder wird vom Wind weggeführt.
\item \emph{Wärmeleitung} bei Kontakt mit kälteren Gegenständen (Sitzfläche, kaltes Wasser).
\item \emph{Verdunstung} von Schweiss: Die Verdampfungswärme wird der Haut entzogen. Bei hoher Umgebungstemperatur und bei Anstrengung ist das der wichtigste Weg.
\item \emph{Atmung}: Die eingeatmete Luft wird in der Lunge erwärmt und mit Wasserdampf gesättigt wieder ausgeatmet.
\end{itemize}
